\documentclass[11pt]{article}

\usepackage[margin=1in]{geometry}
\usepackage{amsmath,amssymb}
\usepackage{graphicx}
\usepackage{booktabs}
\usepackage{array}
\usepackage[colorlinks=true,linkcolor=blue,citecolor=blue,urlcolor=blue]{hyperref}
\usepackage{natbib}
\usepackage{xcolor}
\usepackage{caption}

\newcommand{\code}[1]{\texttt{#1}}

\title{\textbf{Reprogramming a Frozen LLM for Motor Temperature Forecasting:\\
A Multi-Session, Multi-Seed Reproduction and Baseline Study}}

\author{
  Jeevan Merkaji Somanna \\
  \textit{[Affiliation]} \\
  \texttt{jeevanms5355@gmail.com}
}

\date{\today}

\begin{document}

\maketitle

\begin{abstract}
We reproduce a simplified version of Time-LLM \citep{jin2023time}, which forecasts
time series by reprogramming a frozen, pretrained language model rather than
training a forecaster from scratch. Using a GPT-2 (124M) backbone
\citep{radford2019gpt2} on CPU-only hardware, we implement the reprogramming
mechanism and reversible instance normalization (RevIN) \citep{kim2022revin}, and
evaluate against two baselines -- a from-scratch LSTM and a linear DLinear model
\citep{zeng2023dlinear} -- on Kaggle's real-world Electric Motor Temperature
dataset \citep{kagglemotor,kirchgassner2021motor}. Unlike a typical single-run
comparison, we evaluate across \textbf{3 independent real-world recording
sessions} (\code{profile\_id} 70, 62, 44) \(\times\) \textbf{3 random seeds}
(9 runs per method, 27 total), reporting mean \(\pm\) standard deviation rather
than single-point estimates. Our central finding is that \textbf{no single
method dominates across sessions}: the from-scratch LSTM baseline wins on
\code{profile\_id=70} (31\% lower MSE than Time-LLM, replicating an earlier
single-session result), DLinear and Time-LLM are statistically indistinguishable
at a near-zero error floor on \code{profile\_id=62}, and \textbf{Time-LLM wins
outright on \code{profile\_id=44}} (27\% lower MSE than LSTM, 10\% lower than
DLinear). Averaged across all three sessions, Time-LLM achieves the lowest mean
test MSE of the three methods (0.0290 vs.\ LSTM's 0.0337 and DLinear's 0.0348),
while LSTM achieves the lowest mean MAE (0.0581 vs.\ Time-LLM's 0.0793) -- a
split verdict that is itself the main result. This directly complicates the
common assumption, supported by our own earlier single-session pilot, that a
small from-scratch baseline uniformly dominates LLM-reprogramming on
data-abundant, single-sensor tasks, and demonstrates why multi-session
evaluation is necessary before drawing method-level conclusions from
time-series forecasting comparisons -- a single session, however carefully
chosen, can misrepresent the general ranking of methods on a dataset with
heterogeneous recording sessions.
\end{abstract}

\section{Motivation}

Forecasting motor temperature from sensor data is a well-studied problem
usually tackled with purpose-built sequence models (LSTMs, TCNs, Transformers
trained from scratch on the target domain). Time-LLM \citep{jin2023time}
proposes an alternative: keep a large pretrained language model frozen, and
train only a lightweight ``reprogramming'' layer that translates numeric
time-series patches into something resembling the LLM's native
token-embedding space. The appeal is that the frozen LLM contributes general
sequence-reasoning ability learned from massive pretraining, without the cost
of training a large model from scratch or fine-tuning it.

This project asks a practical question relevant to anyone considering this
approach for a real sensor-forecasting task: \textit{on modest hardware, with
a small frozen backbone, does reprogramming actually beat a simple
task-specific baseline you could train in an afternoon?} We answer this
empirically rather than assuming the paper's headline results (obtained with
a 7B-parameter backbone and a full training budget) transfer to a
constrained, CPU-only, GPT-2-scale setting.

A first pass at this question (single session, single seed) found the
from-scratch LSTM baseline winning outright. But a single session/seed
comparison risks over-generalizing from what could be session-specific noise
or an artifact of one random initialization -- a known concern in
time-series forecasting benchmarks, where results can be sensitive to both
the specific recording window and training stochasticity. This work extends
that comparison to \textbf{3 independent real-world sessions \(\times\) 3
seeds per method}, asking: \textit{does the from-scratch baseline's advantage
hold up across multiple independent slices of real data, or was the original
single-session result an artifact of that one session?} The answer, reported
in Section~\ref{sec:results}, is that it was partly an artifact: the LSTM's
advantage replicates on the original session but reverses on another, with
the third session landing all three methods within noise of each other.

\section{Method}

Time-LLM's core mechanism, as implemented here:

\begin{enumerate}
  \item \textbf{Patch embedding.} The input window is split into overlapping
    patches (\code{patch\_len=16}, \code{stride=8}) and linearly projected
    to a small embedding dimension \code{d\_model}.
  \item \textbf{Patch reprogramming.} A multi-head cross-attention layer
    lets each patch embedding attend over a fixed subsample of the frozen
    LLM's own vocabulary embedding matrix (``text prototypes''), producing
    an embedding in the LLM's native space without ever editing the LLM's
    weights.
  \item \textbf{RevIN} \citep{kim2022revin}. Each input window is normalized
    using its own mean/std (not a global training-set statistic), and the
    model's output is denormalized with the same per-window statistics
    before the loss is computed. This is the mechanism used to handle
    distribution shift between windows.
  \item \textbf{Output projection.} The frozen LLM's output over the patch
    positions is flattened and linearly projected to the forecast horizon.
\end{enumerate}

Only the patch embedding, reprogramming layer, RevIN's two affine
parameters, and the output projection are trained. The LLM itself never
updates -- in our runs, 0.1--1.1\% of total parameters were trainable,
consistent with the paper's efficiency claims. Prompt-as-Prefix (PaP), the
paper's optional text-conditioning mechanism, is disabled throughout this
matrix (\code{--no\_prompt}) to prioritize seed/session coverage within a
CPU-only compute budget; it was evaluated separately in an earlier
single-session pilot and roughly quintuples per-sample compute.

\textbf{Baselines.}
\begin{itemize}
  \item \textbf{LSTM} -- 2-layer, 64 hidden units, same RevIN normalization
    and windowing code as Time-LLM, so the comparison isolates the
    forecasting method rather than data-handling differences.
  \item \textbf{DLinear} \citep{zeng2023dlinear} -- a simple linear
    trend/seasonal decomposition model, included as a fast, near-zero-capacity
    reference point.
\end{itemize}

\textbf{Deviations from the original paper}, made for tractability on
CPU-only, no-GPU hardware: a GPT-2 (124M) backbone instead of the paper's
primary Llama-7B; smaller \code{d\_model}/prototype counts and shorter
training schedules than the paper's tuned, per-dataset configurations.

\section{Experimental Setup}

\textbf{Dataset.} Kaggle's Electric Motor Temperature dataset
\citep{kagglemotor,kirchgassner2021motor} (Paderborn University): 185 hours
of PMSM (permanent magnet synchronous motor) test-bench recordings at 2Hz,
69 independent measurement sessions (\code{profile\_id}) concatenated with
no timestamp column. We select 3 sessions to test cross-session robustness:

\begin{center}
\begin{tabular}{@{}lrl@{}}
\toprule
\code{profile\_id} & Rows & Notes \\
\midrule
70 & 25{,}677 & mid-sized session, used in the single-session pilot \\
62 & 25{,}600 & -- \\
44 & 26{,}341 & -- \\
\bottomrule
\end{tabular}
\end{center}

Target: \code{pm} (permanent-magnet/rotor temperature). Window size
\code{seq\_len=240}/\code{pred\_len=60} (2 minutes of history \(\to\) 30
seconds ahead), chosen to avoid the near-trivial persistence regime of
shorter windows at 2Hz sampling.

\textbf{Seeds.} Each (session, method) combination is repeated for seeds
\(\{0, 1, 2\}\), varying model initialization and stochastic training
elements; data splits (chronological 70/10/20 train/val/test) are fixed per
session.

\textbf{Training budget.} DLinear/LSTM: 30 epochs. Time-LLM: 15 epochs.
Batch size 16 throughout. All runs CPU-only (no GPU), frozen GPT-2 (124M)
backbone for Time-LLM.

\textbf{Orchestration.} All 27 runs were driven sequentially by a resumable
orchestration script, with each result checkpointed immediately so progress
survives interruption. Aggregation (mean \(\pm\) std per session/method) is
computed from the raw per-run log.

\section{Results}
\label{sec:results}

All 27/27 runs (3 sessions \(\times\) 3 methods \(\times\) 3 seeds) are
complete.

\subsection{Per-session results}

\begin{table}[h]
\centering
\caption{\code{profile\_id=70}}
\begin{tabular}{@{}lccc@{}}
\toprule
Method & Test MSE & Test MAE & Train time (mean) \\
\midrule
DLinear & $0.0350 \pm 0.0073$ & $0.1330 \pm 0.0247$ & 0.05h \\
\textbf{LSTM} & $\mathbf{0.0150 \pm 0.0007}$ & $\mathbf{0.0629 \pm 0.0041}$ & 1.58h \\
Time-LLM & $0.0247 \pm 0.0058$ & $0.0985 \pm 0.0218$ & 8.13h \\
\bottomrule
\end{tabular}
\end{table}

\begin{table}[h]
\centering
\caption{\code{profile\_id=62}}
\begin{tabular}{@{}lccc@{}}
\toprule
Method & Test MSE & Test MAE & Train time (mean) \\
\midrule
\textbf{DLinear} & $\mathbf{0.0003 \pm 0.0000}$ & $\mathbf{0.0135 \pm 0.0001}$ & 0.04h \\
LSTM & $0.0007 \pm 0.0005$ & $0.0173 \pm 0.0049$ & 1.47h \\
Time-LLM & $0.0004 \pm 0.0002$ & $0.0157 \pm 0.0029$ & 3.70h \\
\bottomrule
\end{tabular}
\end{table}

\begin{table}[h]
\centering
\caption{\code{profile\_id=44}}
\begin{tabular}{@{}lccc@{}}
\toprule
Method & Test MSE & Test MAE & Train time (mean) \\
\midrule
DLinear & $0.0691 \pm 0.0208$ & $0.0985 \pm 0.0049$ & 0.05h \\
LSTM & $0.0853 \pm 0.0086$ & $0.0940 \pm 0.0020$ & 0.53h \\
\textbf{Time-LLM} & $\mathbf{0.0619 \pm 0.0019}$ & $0.1236 \pm 0.0247$ & 6.58h \\
\bottomrule
\end{tabular}
\end{table}

\textbf{Each method wins exactly one of the three sessions on test MSE.} On
\code{profile\_id=70}, LSTM beats Time-LLM by 31\% MSE -- a clean replication
of the single-session pilot. On \code{profile\_id=62}, all three methods sit
within a very narrow, near-zero error band (this session's \code{pm} target
appears to have an unusually low dynamic range/noise floor, making the
forecasting task close to saturated for all three methods regardless of
capacity); DLinear is nominally lowest but the seed-to-seed spread mostly
overlaps with Time-LLM's. On \code{profile\_id=44}, \textbf{Time-LLM wins
outright}: 27\% lower MSE than LSTM and 10\% lower than DLinear, with the
smallest seed-to-seed std of the three methods on this session
($\pm 0.0019$ MSE) -- the most confidently-won session in the whole matrix.

\subsection{Cross-session summary}

\begin{table}[h]
\centering
\caption{Mean of per-session means, unweighted across sessions.}
\begin{tabular}{@{}lccc@{}}
\toprule
Method & Mean Test MSE & Mean Test MAE & Sessions won (MSE) \\
\midrule
DLinear & 0.0348 & 0.0817 & 1/3 (\code{p62}) \\
LSTM & 0.0337 & \textbf{0.0581} & 1/3 (\code{p70}) \\
\textbf{Time-LLM} & \textbf{0.0290} & 0.0793 & 1/3 (\code{p44}) \\
\bottomrule
\end{tabular}
\end{table}

(See Section~\ref{sec:discussion} for a caveat on why this specific
unweighted aggregation should be read cautiously despite Time-LLM's lead
here.)

\section{Discussion}
\label{sec:discussion}

\textbf{No method dominates across sessions -- the split verdict is the
finding.} Each of the three methods wins exactly one session on test MSE,
and the cross-session mean and the cross-session ``sessions won'' count
don't even agree on a single winner by every metric: Time-LLM has the
lowest mean MSE, but LSTM has the lowest mean MAE, and all three are tied
1-1-1 on session wins. This is the central methodological point of this
study: a single-session comparison, however carefully run (the original
pilot used a full 15/30-epoch budget, proper RevIN normalization, and a
fair from-scratch baseline -- nothing about the original comparison was
sloppy) can still produce a method-ranking claim that reverses on a second,
equally legitimate session of the same underlying real-world dataset.
Multi-session evaluation is not a nice-to-have for honest claims about
which forecasting approach ``wins'' on a sensor-data problem -- per-session
dynamics (noise floor, signal range, operating regime) can matter as much
as, or more than, the modeling choice itself.

\textbf{The cross-session mean MSE should be read with a specific caveat.}
\code{profile\_id=62}'s errors are roughly two orders of magnitude smaller
than the other two sessions' (0.0003--0.0007 vs.\ 0.015--0.09), so an
unweighted mean across sessions is dominated by whichever method happens to
be marginally better on the two ``harder'' sessions (\code{p70}, \code{p44})
rather than reflecting a genuinely comparable average. Time-LLM's
cross-session MSE lead is real but should be understood as ``wins the
harder session (\code{p44}) outright, loses the other harder session
(\code{p70}), ties on the easy one (\code{p62})'' rather than ``wins on
average'' in any deeper sense -- we report the unweighted mean because it
is the simplest defensible aggregation given only 3 sessions, not because
we think it is the last word on ranking.

\textbf{Why might Time-LLM win specifically on \code{profile\_id=44}?} We
do not have a confirmed mechanistic answer -- this would require inspecting
each session's underlying operating conditions (load profile, thermal
transient vs.\ steady-state behavior) in the source PMSM dataset, which is
out of scope here. What we can say: \code{profile\_id=44}'s std across
seeds for Time-LLM ($\pm 0.0019$) is tighter than LSTM's ($\pm 0.0086$) and
DLinear's ($\pm 0.0208$) on the same session, meaning Time-LLM's advantage
there is not an artifact of a single lucky seed. This is consistent with
(though not proof of) the reprogrammed LLM's frozen pretrained backbone
providing more useful inductive bias on sessions with more
complex/nonstationary dynamics than a small from-scratch model can pick up
in 15--30 epochs on one session's data alone.

\textbf{Compute cost remains a first-class, session-independent result.}
Across all sessions, Time-LLM training took roughly 4--16$\times$ longer
than the LSTM baseline (1.5--8.1h vs.\ 0.5--1.6h) for accuracy that is
session-dependently better, worse, or tied. This tradeoff is directly
decision-relevant regardless of which method wins on accuracy for a given
session, and matters even more once the result is ``it depends which
session you're on'' rather than a clean win for either side.

\section{Limitations and Future Work}

\begin{itemize}
  \item \textbf{Backbone scale} is the most likely lever left untested:
    repeating these experiments with a 7B-parameter backbone (as in the
    original paper) on a GPU would directly test whether the session-split
    result changes with more model capacity.
  \item \textbf{Prompt-as-Prefix disabled} in this matrix for CPU-time
    reasons; a full-rigor PaP comparison (matched epochs, richer
    natural-language prompt) is a natural extension.
  \item \textbf{3 of 69 available sessions} were used; broader coverage (more
    sessions, or a formal stratified sample across operating conditions)
    would further strengthen the cross-session claim.
  \item \textbf{Time-LLM epoch budget (15) vs.\ LSTM/DLinear (30)} --
    asymmetric budgets are a possible confound worth noting, kept for
    CPU-time reasons.
\end{itemize}

\section{Reproducibility}

All code, exact hyperparameters, per-run logs, and the raw results file for
every run reported here are available at:
\begin{center}
\url{https://github.com/Jeevanmerkaji/Time-LLM}
\end{center}

\bibliographystyle{plainnat}
\bibliography{references}

\end{document}
